% Conclusion (draft)
\label{sec:conclusion}
This paper presented a leakage-controlled evaluation of football match
outcome forecasting in which statistical models, tree-based models,
and a domain-enhanced large language model are compared against
de-vigged closing odds under an identical, fully specified protocol.
On 1,104 held-out matches, no learned method exceeded market
accuracy, and the LLM matched market-level probability calibration
at negligible cost. Feature importance and ablations show that market
signals dominate, while referee statistics add nothing measurable
and xG adds a small consistent gain. The principal positive result
is that uncertainty-driven risk control works: an uncertainty index
and a corrected scenario-complexity score stratify held-out matches
monotonically, a no-bet policy on high-uncertainty matches improves
risk-adjusted returns, and draw outcomes are systematically
under-predicted, with class weighting providing a controllable
accuracy--balance trade-off.

These results frame a realistic division of labor for
decision-support systems in sports betting: treat the market as the
prior, use models for calibrated probability estimates and
explainability, and rely on uncertainty signals---not claimed
accuracy---for risk control. Forecast difficulty---the ex-ante
probability that the system is wrong---is the quantity a
decision-support system can realistically deliver in an efficient
market, and abstention on the basis of that quantity is the
decision-level counterpart of selective prediction. To support
reproducibility, the full
data-processing pipeline and experiment code are available at
\url{https://github.com/wenjin06/football-outcome-forecasting-benchmark};
no proprietary data are required beyond public sources.
